\section{Implementation, practical use and testing practice}
\label{sec:implementation}
% P:implementation-01
Every iterative calculation runs natively in C++: association and density closure, equilibrium solves, phase searches, continuation and parameter fitting.
The Python interface declares inputs and reads results; it supplies no callback to a native iteration and performs no thermodynamic arithmetic.
A caller constructs a model once from a parameter record.
Construction validates the species, site topology and model selection, admits the records of Section~\ref{sec:references}, and prepares state-independent quantities; support and thermal intervals select reusable recordings of the shared energy equations.
A state request then gives temperature, composition and exactly one of density or pressure; a pressure request resolves the declared density branch once and reuses that state for every property and derivative asked of it.
An equilibrium request is a declaration of phases, supports, pinned amounts, free intensive variables, feed and reactions, compiled once into the rows of Section~\ref{sec:equilibrium}.
Figure~\ref{fig:workflow} shows how these requests share one calculation.

% P:implementation-02
Listing~\ref{lst:example} declares a methane/ethane mixture with the parameters of Gross and Sadowski~\citep{Gross2001}, solves a bubble point at 200~K and requests the exact first-order response of the bubble pressure to temperature.
The same request with feed or parameter directions, or with phase compositions as observables, returns the other columns of $\mathbf{y}_\theta$ for that state.

\noindent\begin{minipage}{\linewidth}
\begin{lstlisting}[caption={A bubble point and its exact temperature derivative.},label={lst:example}]
import epcsaft
from epcsaft import equilibrium as eq

record = "data/packets/gross-2001-methane-ethane/3/parameters/parameters.json"
mix = epcsaft.Mixture(epcsaft.Parameters.from_json(record, components=["methane", "ethane"]))
feed = eq.MoleFractions({"methane": 0.3, "ethane": 0.7})
problem = eq.bubble_point(liquid=eq.Phase("L", kind="liquid"), T=200.0, feed=feed)
bubble = eq.solve_equilibrium(mix, problem)          # pressure is free
print(bubble.pressure, bubble.phase_mole_fractions(1))   # Pa; incipient vapor

dT = eq.SolvedStateActionDirection(1.0, 0.0, [0.0, 0.0], [ ])  # unit temperature direction
P = eq.SolvedStateObservable(eq.SolvedStateObservableKind.Pressure)
batch = eq.solved_state_actions(eq.compile_problem(mix, problem), bubble,
                                eq.SolvedStateActionRequest([P], [dT]))
print(batch.results[0].action, batch.results[0].status)       # Pa/K
\end{lstlisting}
\end{minipage}

% P:implementation-03
Equation-oriented process models reach the same calculation through a compiled library that implements external functions of the AMPL Solver Library.
A model written in Pyomo~\citep{hartPyomoModelingSolving2011}, for example within the IDAES process modeling framework~\citep{leeIDAESProcessModeling2021}, calls these functions and receives values with exact first and second derivatives.
The functions verified here return the pressure of a neutral fluid with its gradient and Hessian, and the phase compositions and amounts of a declared fixed-temperature, fixed-pressure, nonreactive two-liquid electrolyte split with their first and second feed derivatives.
It solves each call with the same native driver and adds no phase search.

% P:implementation-04
The native code depends on CppAD~20250000.3~\citep{bellCppADPackageAlgorithmic2025} for algorithmic differentiation, Ipopt~3.11.9~\citep{Wachter2006} with the MUMPS sparse linear solver~\citep{amestoyFullyAsynchronousMultifrontal2001,amestoyPerformanceScalabilityBlock2019} for equilibrium and minimax problems, Ceres~2.2~\citep{agarwalCeresSolver2023} for least squares, Eigen~3.4~\citep{guennebaudEigenV32010} for dense linear algebra, and HiGHS~1.15.1~\citep{huangfuParallelizingDualRevised2018} for the conserved-inventory linear programs.
The Python interface is built with nanobind~2.12.0~\citep{jakobNanobindTinyEfficient2022}.


\subsection{Testing and verification practice}
% P:implementation-05
Checks follow the dependency of a request: contribution equations, implicit closures, equilibrium balances and phase identity, derivative actions, then observation composition and process transport.
This helps distinguish an accurate property evaluation from a wrong root or an inaccurate physical model.


% P:implementation-06
The comparisons include 704 branch-specific density requests, 36 pressure-specified caloric actions and 3\,167 electrolyte observation columns, with reference values, identities, independent codes and refinement chosen by quantity.
Table~\ref{tab:verification} states representative numerical bounds.
Counts describe checked requests or scalar comparisons, not independent experiments.


% P:implementation-07
Practical use starts with a parameter record and specification, followed by inspection of achieved residuals, unavailable causes and source domains.
For external functions, the process model owns its balances, constraints and optimization, while the package owns the supplied properties or declared equilibrium and their derivatives.
An external flash may therefore assemble its own equilibrium equations from the package's properties, as illustrated in Section~\ref{sec:applications}.
